%% ============================================================================
%% DRAFT for Jared, 2026-09-21 (numbers refreshed from the v1.3.1 rerun,
%% multislice n_starts = 10, same evening). A main-text subsection: when to use
%% which method. Sources:
%% tab_decision_tree.tex, tab_metrics_wide.tex, tab_app_selection_rule.tex,
%% tab_app_empirical_stability.tex). Adopt, cut, or rewrite; nothing here is
%% in the Overleaf project yet. Table reference assumes
%% \input{tables/tab_decision_tree.tex} with \label{tab:decision_tree}.
%% ============================================================================

\subsection{When to Use DynMux, Multislice, or Label Matching}
\label{sec:when_to_use}

The three trackers in the \texttt{dynamicmultiplex} package differ in what
they carry across layers. DynMux couples \emph{communities}: a community at
$t$ is linked to a community at $t'$ when their memberships overlap, and no
node is tied to itself. Multislice modularity couples \emph{nodes}: node $i$
at $t$ is tied to node $i$ at $t+1$ with weight $\omega$, so the objective
rewards every node that stays put. Hungarian label matching couples nothing;
it detects each layer on its own and aligns labels afterwards. Which coupling
is right depends on whether the thing that persists over time is the node or
the community, and Table~\ref{tab:decision_tree} collects the evidence from
every simulated scenario in the paper.

\begin{table}[t]
\centering\small
\caption{Joint NMI of the tracked partition by scenario (means over
replicates; Table 2 regimes use 30 replicates per configuration, the
mechanism and lineage regimes 10 and 30). Multislice uses adjacent-layer
links except where the generator supplies period links. Margin is the gap
between the best method and the runner-up.}
\label{tab:decision_tree}
\input{tables/tab_decision_tree.tex}
\end{table}

\paragraph{Use DynMux when the node set or the community set changes.}
DynMux is the most accurate tracker when communities are born and die
(joint NMI 0.37 against 0.31 for multislice and 0.28 for label matching),
when nodes enter and leave the network (0.36 against 0.28; DynMux is the more
accurate method in 92\% of series), and when structure recurs on a period
the analyst cannot supply as a link (0.28 against 0.10). In each case the
identity tie has nothing to attach to, or attaches the wrong thing: an
absent node has no copy in the next layer, a dead community's members are
dragged toward its old label, and adjacent-layer ties cannot see a
partition that returns after a gap. Community-level coupling ignores all
three problems because it never asks whether node $i$ is the same node.
%% [CLAUDE 2026-09-21, v1.3.1] Abrupt rewiring: with 10 multislice restarts the Table 2 tie holds (DynMux 0.490 vs multislice 0.486), so the restart-noise explanation for the earlier omega-sweep discrepancy is ruled out; the omega sweep still puts multislice at 0.41 on this regime, and my working hypothesis is that sim/05's 72-config grid is not the same grid as sim/01 (to be checked by diffing the two config tables; not a manuscript blocker because the omega table is only used to show insensitivity to omega). Safe sentence for the text: "at an abrupt change-point the two are indistinguishable on joint NMI (0.49 each), and DynMux has the lowest K error of any method (1.36 against 1.37 and 1.40 for the two multislice variants)". Do not list abrupt rewiring as a DynMux win.

\paragraph{Use multislice when nodes persist and the signal is weak.}
When most nodes keep their community from one layer to the next, the
identity tie is the right prior and it also pools evidence across layers.
Multislice wins by wide margins under gradual switching (0.53 against 0.28),
under a stable core with a rotating periphery (0.86 against 0.46), and on
long series with many small communities (0.69 against 0.23), where DynMux's
second stage over-splits ($K$ error 5.3 against 1.5). It also wins in every
lineage regime of Appendix~X (splits, merges, shrinking and nested
communities), because its supra-graph carries a community through the event
as one object while DynMux's community graph fragments it. When the period
of a recurring structure is known, multislice with period links matches or
beats DynMux (0.43 against 0.28). None of these results depend on $\omega$:
across $\omega \in \{0.25, 0.5, 1, 2, 4\}$ the multislice numbers move by
at most 0.012 (Appendix~X).

\paragraph{Use label matching only to preserve per-layer partitions.}
Hungarian matching is never the most accurate method in any scenario. Its
use case is procedural: when the analyst needs each layer's partition exactly
as a single-layer algorithm returns it, with labels aligned afterwards and no
information shared across layers.

\paragraph{When the regime is unknown, compare stability scores.}
Fitting both DynMux and multislice and keeping the one with the higher
stability score (Section~4) selects the more accurate method in 82\% of
simulated series across the four Table~2 regimes, and yields a mean joint
NMI of 0.47 against 0.35 for always choosing DynMux, 0.44 for always
choosing multislice, and 0.48 for an oracle that knows the answer
(Appendix~X). On the four empirical networks the rule prefers multislice in
every case (Appendix~X): its stability scores are 0.97, 0.80, 0.81 and 0.59
on the alliance, defense cooperation, IGO and trade networks against 0.85,
0.51, 0.72 and 0.12 for DynMux. The two methods recover the coded orders
about equally well (Section~5), so the difference is in how confidently the
node-level assignments can be read, and Section~5 reports both.
%% [CLAUDE 2026-09-21] v1.3.1 changed only the DCA (0.75 -> 0.80) and IGO (0.80 -> 0.81) multislice scores; DynMux scores are unchanged. The last two sentences are the honest version of the empirical stability result; they weaken the "DynMux is the method for the empirical networks" framing in Sec. 5 and the conclusion. If you would rather keep Sec. 5 as is, this paragraph still has to state the empirical stability numbers somewhere, because tab_app_empirical_stability.tex is in the appendix.

%% [CLAUDE 2026-09-21] Overlap: one sentence somewhere in Sec. 3.1, not here: "An overlap-coefficient coupling is also available; it never beats Jaccard on membership accuracy and is worse when communities split (-0.05) or merge (-0.06), so the main text uses Jaccard throughout (Appendix X)."
